\chapter{Usage}

This chapter describes how \cf{} is started and how it is used in practice. The
system exposes the same reasoning engine through two interfaces: a command-line
REPL, useful for development and testing, and a Telegram bot, which is
the intended everyday interface ---usable at the pub, not in front of a computer.
Both interfaces drive the same dialogue state machine and the same knowledge
base, so a consultation conducted from the terminal and one conducted from the
bot ask the very same questions and reach the very same recommendations. The
chapter first explains how to start each interface, then enumerates the available
commands, and finally walks through a number of worked sessions ---a full
consultation, a meal-pairing consultation, a comparison of the two reasoning
backends, the \emph{Help} and \emph{Why} facilities, and the revision and
cancellation of answers--- each illustrated with a faithful transcript.


\section{Starting the system}

The reasoning engine runs from the terminal with a CLIPS~6.4.x interpreter. The
project does not assume a system-wide CLIPS: the helper script
\code{scripts/build-clips.sh} builds CLIPS~6.4.2 locally into \code{\~{}/.local/bin},
after which the engine is launched by loading its start-up file
\code{BeerEX.clp}:

\begin{lstlisting}[language=bash,numbers=none]
bash scripts/build-clips.sh    # builds CLIPS 6.4.2 into ~/.local/bin
clips -f BeerEX.clp
\end{lstlisting}

The start-up file loads the core templates and globals, the fuzzy backend, the
per-style numeric ranges and profiles, the dialogue state machine and, finally,
the question and knowledge rule files; it then performs a \code{(reset)} and
\code{(run)} and hands control to the dialogue loop, which prints the first
question and waits for the user's answer. From then on the program asks the
questions one at a time, accepting either one of the displayed options or one of
the \emph{help}, \emph{why}, \emph{prev} (back) and \emph{cancel} control words.

The Telegram bot, by contrast, needs no system CLIPS: the \code{clipspy} binding
bundles its own interpreter. The bot is built on \emph{python-telegram-bot}~v21
and keeps one isolated CLIPS environment per chat, so that concurrent users never
share working memory. With a token obtained from \emph{@BotFather}, the bot is
started either through Docker or directly with Python~3.12:

\begin{lstlisting}[language=bash,numbers=none]
cp .env.example .env                   # put your @BotFather TOKEN in .env
docker compose up --build              # recommended

# or, for a local run:
git submodule update --init            # fetch the fuzzy library
python3.12 -m venv .venv && . .venv/bin/activate
pip install -r requirements.txt
export TOKEN=<your-telegram-bot-token>
python BeerEX.bot.py
\end{lstlisting}

Once the process is polling, the user opens the chat and types \code{/start}. The
bot greets the user by name and prints a short presentation, after which a
consultation can be started with \code{/new}. The options for each question are
offered as a reply keyboard, together with the \emph{Help}, \emph{Why},
\emph{Previous} and \emph{Cancel} buttons where applicable; when a recommendation
has been produced, the keyboard offers instead the \emph{Previous},
\emph{Restart} and \emph{Cancel} buttons.


\section{Commands}

The bot recognises a small set of slash commands, which control the session as a
whole, and a set of in-dialogue buttons, which drive the question/answer
interaction. The slash commands are the following:

\begin{itemize}
  \item \code{/start} --- greets the user and prints a short description of how
        the system works;
  \item \code{/new} --- starts a new consultation, resetting the working memory
        and asking the first question;
  \item \code{/cf} --- selects the Certainty-Factors reasoning backend (the
        classic MYCIN-style evidence combination);
  \item \code{/fuzzy} --- selects the fuzzy-logic reasoning backend, which scores
        every style with data-grounded membership functions;
  \item \code{/en} --- sets the output language to English (the base language);
  \item \code{/it} --- sets the output language to the Italian overlay, which
        falls back to English wherever a string is not yet translated.
\end{itemize}

The \code{/cf}, \code{/fuzzy}, \code{/en} and \code{/it} commands do not affect
the consultation already in progress: they record the requested backend or
language on the session and take effect at the next \code{/new}. English is
always the base; the commands themselves are control tokens and are never
localised, only the questions, recommendations and explanations that the user
reads are.

During a consultation the following in-dialogue controls are available, exposed
as keyboard buttons (and as the corresponding control words \code{help},
\code{why}, \code{prev}, \code{restart} and \code{cancel} on the command line):

\begin{itemize}
  \item \emph{Help} --- shows an extended description of the current question's
        options; offered only when the question carries a help text;
  \item \emph{Why} --- explains why the current question is being asked; offered
        only when the question carries a rationale;
  \item \emph{Previous} --- goes back to the preceding question, retracting the
        last answer so that it can be revised; offered only once at least one
        question has been answered;
  \item \emph{Restart} --- abandons the current line of questioning and begins a
        fresh consultation; offered on the final recommendation screen;
  \item \emph{Cancel} --- closes the consultation and clears the working memory.
\end{itemize}


\subsection{A worked consultation}

A full consultation begins with a block of profiling questions ---asked in random
order, since they are independent of one another--- covering the drinker's
profile (age, regular-drinker status, smoking habit, season) and their
personal preferences (colour, alcohol, carbonation, flavour). Once these are
known, the dialogue switches to a depth-first chain of meal questions. The
following command-line transcript drives the profiling questions to the point
where a recommendation is produced; the questions and their options are exactly
those defined in the knowledge base.

\begin{lstlisting}[numbers=none]
How old are you?
1.18-23 2.24-29 3.>=30 4.why 5.cancel
Answer: 3

Are you a regular beer drinker?
1.yes 2.no 3.why 4.cancel 5.prev
Answer: yes

Do you usually smoke while drinking?
1.yes 2.no 3.why 4.cancel 5.prev
Answer: no

Is it autumn, spring, summer or winter?
1.autumn 2.spring 3.summer 4.winter 5.why 6.cancel 7.prev
Answer: winter

Do you generally prefer pale, amber, brown or dark beer?
1.pale 2.amber 3.brown 4.dark 5.don't mind 6.cancel 7.prev
Answer: dark

Do you generally prefer to drink low, mild, high or very high alcoholic drinks?
1.low 2.mild 3.high 4.very high 5.don't mind 6.cancel 7.prev
Answer: high

Do you generally prefer to drink low, medium or high carbonated drinks?
1.low 2.medium 3.high 4.don't mind 5.cancel 6.prev
Answer: medium

What kind of flavor do you generally prefer?
1.clean 2.sweet 3.bitter 4.roasty 5.fruity 6.spicy 7.sour 8.don't mind 9.cancel 10.prev
Answer: roasty
\end{lstlisting}

At this point the drinker's profile is complete and the dialogue moves on to the
meal. For the purposes of this example the drinker declares that no particular
dish is involved; the engine combines the accumulated evidence and prints the
recommendation. With the fuzzy backend active, the result block names each style,
its serving hint and its graded certainty, and is followed by the narrative
explanation of the inferred profile together with a per-beer rationale:

\begin{lstlisting}[numbers=none]
Is the main meal pizza, entrée, cheese or dessert?
1.pizza 2.entrée 3.cheese 4.dessert 5.other 6.cancel 7.prev
Answer: other

✅ Done. I have selected these beer styles for you...

🍺 German-Style Schwarzbier — 🌡 45-50°F · 🥃 Flute    with certainty 71%
🍺 German-Style Dunkel — 🌡 45-50°F · 🥃 Vase          with certainty 66%
🍺 American Stout — 🌡 50-55°F · 🥃 Tulip              with certainty 63%

... for the following reasons:

Your palate could be very demanding: I bet you would try out something
peculiar. Frosty days are perfect to taste something sweet or roasty, that
strongly change your palate. Also, high alcoholic beers can heat you up.
You prefer a dark beer. You chose a high alcoholic beer. It's a medium
carbonated beer. You desire a beer with a roasty flavor.

🔎 Why these beers:
🍺 German-Style Schwarzbier: dark colour (μ1.00), noticeable alcohol
(μ0.85), medium carbonation, dark-roasty flavour.
\end{lstlisting}

The narrative explanation is assembled, in order, from the scenario, preference,
main-meal and specific-meal rationales accumulated during the consultation, so
the user can always trace the recommendation back to the answers that produced
it. Under the fuzzy backend a final \emph{Why these beers} paragraph reports, for
each recommended style, the data-grounded membership values that earned it its
place in the ranking.


\subsection{Meal pairing}

When the drinker names a dish, the dialogue descends a meal-specific chain of
questions and the recommendation becomes pairing-driven: the authoritative
CraftBeer food pairings stored in the knowledge base are consulted, and styles
whose pairing matches one of the foods the user named are rewarded. The
transcript below assumes the profiling questions have already been answered as in
the previous session, and picks up at the main-meal question, following the
\emph{pizza} branch to a specific topping:

\begin{lstlisting}[numbers=none]
Is the main meal pizza, entrée, cheese or dessert?
1.pizza 2.entrée 3.cheese 4.dessert 5.other 6.cancel 7.prev
Answer: pizza

Is the pizza topping classic, meat, vegetables, cheese or other?
1.classic 2.meat 3.vegetables 4.cheese 5.other 6.why 7.cancel 8.prev
Answer: meat

Is the meat spicy?
1.yes 2.no 3.why 4.cancel 5.prev
Answer: no

✅ Done. I have selected these beer styles for you...

🍺 American IPA — 🌡 50-55°F · 🥃 Tulip               with certainty 78%
🍺 English-Style IPA — 🌡 45-50°F · 🥃 Nonic Pint     with certainty 74%
🍺 Imperial IPA — 🌡 50-55°F · 🥃 Tulip               with certainty 70%

... for the following reasons:

You're eating a pizza with meat: the best way to enjoy it is pair that with
beers with a bitter flavor, because of the bitterness of the hops that
pleasantly cleanses the mouth.
\end{lstlisting}

The non-spicy meat topping makes the knowledge base assert a strong preference
for a \emph{hoppy-bitter} flavour, which the fuzzy backend translates into a
desire for high bitterness on the IBU axis; the hop-forward styles
therefore rise to the top of the ranking. Had the meat been declared spicy, the
opposite evidence would have been asserted ---a malty-sweet flavour to neutralise
the capsaicin--- and the ranking would have shifted towards the malt-forward
styles instead.


\subsection{Switching the backend}

The same profile can be evaluated under either reasoning backend, and the two
produce visibly different rankings. The Certainty-Factors backend combines the
discrete evidence asserted by the knowledge base in the MYCIN style: a style is
scored by the certainty factors that vote for it, which yields a coarse ranking
in which several styles frequently share the same certainty and so tie. The fuzzy
backend, by contrast, scores every style by a normalised weighted mean of
data-grounded membership values ---colour against SRM, alcohol against
ABV, bitterness against IBU--- so that ties are broken by how
well each style actually fits the desired profile, and the certainties come out
as graded percentages rather than as a handful of repeated values.

Selecting the backend is a matter of one slash command, which takes effect at the
next \code{/new}:

\begin{lstlisting}[numbers=none]
/cf
Backend set to cf. Press /new to start.
\end{lstlisting}

Running the same answers as in the worked consultation above under the two
backends illustrates the difference. Under \code{/cf} the top of the ranking
tends to be a block of styles sharing the certainty contributed by the dominant
rule, for example:

\begin{lstlisting}[numbers=none]
✅ Done. I have selected these beer styles for you...

🍺 German-Style Schwarzbier  with certainty 50%
🍺 German-Style Dunkel        with certainty 50%
🍺 American Stout             with certainty 50%
\end{lstlisting}

Under \code{/fuzzy} the same three styles are separated, because the membership
scores grade how closely each one matches the requested dark colour, noticeable
alcohol and roasty flavour, and the serving hint and per-beer rationale are added:

\begin{lstlisting}[numbers=none]
✅ Done. I have selected these beer styles for you...

🍺 German-Style Schwarzbier — 🌡 45-50°F · 🥃 Flute    with certainty 71%
🍺 German-Style Dunkel — 🌡 45-50°F · 🥃 Vase          with certainty 66%
🍺 American Stout — 🌡 50-55°F · 🥃 Tulip              with certainty 63%
\end{lstlisting}

The Certainty-Factors ranking is therefore useful as a quick, transparent view of
which rules fired, while the fuzzy ranking is preferable whenever a finer, fully
ordered recommendation is wanted. Because both backends feed the same final
formatting stage, the explanation and the result layout are identical in
structure; only the certainties, the ordering and the presence of the fuzzy
rationale differ.


\subsection{Help and Why}

Every question may carry two optional pieces of supporting text: a \emph{help}
text, which explains the meaning of the options, and a \emph{why} text, which
justifies the question. When present, they are offered as the \emph{Help} and
\emph{Why} buttons (or the \code{help} and \code{why} control words on the
command line) and do not advance the dialogue: after the text is shown, the same
question is asked again.

The \emph{Why} facility is illustrated on the age question. Requesting \code{why}
prints the rationale and re-poses the question:

\begin{lstlisting}[numbers=none]
How old are you?
1.18-23 2.24-29 3.>=30 4.why 5.cancel
Answer: why

Younger people tend to prefer less intense beers.

How old are you?
1.18-23 2.24-29 3.>=30 4.why 5.cancel
Answer: 1
\end{lstlisting}

The \emph{Help} facility is illustrated on the cheese-style question, one of the
few questions rich enough to carry a help text describing each option in detail.
Requesting \code{help} prints the extended description of the cheese families
before the question is asked again:

\begin{lstlisting}[numbers=none]
Is the cheese style fresh, semi-soft, firm/hard, blue, natural rind
or washed rind?
1.fresh 2.semi-soft 3.firm/hard 4.blue 5.natural rind 6.washed rind
7.don't know 8.help 9.why 10.cancel 11.prev
Answer: help

🧀 Fresh cheeses have not been aged, or are very slightly cured. These
cheeses have a high moisture content and are usually mild and have a very
creamy taste and soft texture.
🧀 Semi-soft cheeses have a smooth, generally, creamy interior with little
or no rind. [ ... ]
\end{lstlisting}

In the Telegram interface the help text is rendered with Markdown; where it
embeds an image link, the bot turns the link into an inline button that sends the
corresponding picture when pressed.


\subsection{Going back and cancelling}

The \emph{Previous} control lets the user revise an answer. Going back retracts
the answer given to the current question and, with it, any conclusion derived
from that answer: the dialogue state machine removes the fact asserted for the
question and re-fires the affected rules, so that working memory remains
consistent with the answers actually in force. The transcript below shows a
drinker who answered \emph{pale} to the colour question, went back, and changed
the answer to \emph{dark}:

\begin{lstlisting}[numbers=none]
Do you generally prefer pale, amber, brown or dark beer?
1.pale 2.amber 3.brown 4.dark 5.don't mind 6.cancel 7.prev
Answer: pale

Do you generally prefer to drink low, mild, high or very high alcoholic drinks?
1.low 2.mild 3.high 4.very high 5.don't mind 6.cancel 7.prev
Answer: prev

Do you generally prefer pale, amber, brown or dark beer?
1.pale 2.amber 3.brown 4.dark 5.don't mind 6.cancel 7.prev
Answer: dark
\end{lstlisting}

Because the colour evidence asserted for \emph{pale} is retracted when the user
goes back, the subsequent recommendation reflects only the corrected \emph{dark}
preference, with no residual trace of the abandoned answer. The same mechanism
applies even on the final recommendation screen: going back from there retracts
the last meal answer and re-poses it, so the user can explore alternative meal
choices without restarting.

Finally, the \emph{Cancel} control closes the consultation. On the command line
it terminates the interpreter; in the bot it clears the chat's working memory,
removes the keyboard and bids the user farewell:

\begin{lstlisting}[numbers=none]
Answer: cancel
Bye! I hope we can talk again some day.
\end{lstlisting}

After cancelling, a fresh consultation can always be started again with
\code{/new} (or by relaunching the command-line interpreter).
